\subsection{Local Cutoff}
\label{app:manual-local}\label{app:analytic-local}

Write \(x=e^{-\ell}\). We use the single influence cutoff
\begin{equation}
                 a_{\rm cut}=\frac{1}{1+(8/3)x}.
 \label{eq:rational-local-cutoff}
\end{equation}
We need this cutoff only for \(\ell\ge7/2\). On this half-line
\(x<1/32\): indeed \(e^3>20\) and
\(e^{1/2}>1+1/2+1/8=13/8\) give \(e^{7/2}>32\).
The proof retains the leading entropy logarithms until they cancel,
then bounds the three remaining errors uniformly.

Put \(r=(1-x^2)/(1+x^2)\), and let
\(h_b(p)=-p\log p-(1-p)\log(1-p)\) denote binary entropy in
probability coordinates. The following temporary quantities expose
the cancellation between the two entropies:
\begin{align*}
 \kappa_{\rm loc}&=\frac4{3+8x},&
 \eta_{\rm loc}&=\frac{3x}{4(1+x^2)},\\
 p&=\frac{x^2}{1+x^2},&
 q_{\rm loc}&=\kappa_{\rm loc}x(1+\eta_{\rm loc}),\\
 A_{\rm loc}&=\frac1{2x^2}+\frac3{4x}+1+\frac{x^2}{2}.
\end{align*}
Thus \(p=(1-r)/2\), \(q_{\rm loc}=(1-r a_{\rm cut})/2\). The local margin
\((1-r^2)h(r a_{\rm cut})-(1-a_{\rm cut}r^2)h(r)\), divided by the positive factor
\(4\kappa_{\rm loc} x^3/(1+x^2)^2\), equals
\[
 \Lambda_{\rm loc}=\frac{h_b(q_{\rm loc})}{\kappa_{\rm loc}x}
 -A_{\rm loc}h_b(p).
\]

\paragraph{Cancellation of the leading logarithms.}
Both entropy estimates follow from one pair of logarithm bounds:
\[
 \frac{2t}{2+t}\le\log(1+t)\le t-\frac{t^2}{2(1+t)}
 \qquad(t\ge0).
\]
These are the midpoint and trapezoid bounds for the integral of the
convex function \(s\mapsto(1+ts)^{-1}\) on \([0,1]\), multiplied by \(t\).
The lower bound, with \(t=q_{\rm loc}/(1-q_{\rm loc})\), gives
\begin{align*}
 h_b(q_{\rm loc})\ge{}&q_{\rm loc}\log(1/q_{\rm loc})+q_{\rm loc}\\
 &-\frac{q_{\rm loc}^2}{2-q_{\rm loc}}.
\end{align*}
The upper bound gives
\((1+\eta_{\rm loc})\log(1+\eta_{\rm loc})\le\eta_{\rm loc}+\eta_{\rm loc}^2/2\).
Consequently
\begin{align*}
 \frac{h_b(q_{\rm loc})}{\kappa_{\rm loc}x}\ge{}&
 (1+\eta_{\rm loc})(\ell-\log\kappa_{\rm loc})+1\\
 &-\frac{\eta_{\rm loc}^2}{2}
 -\frac{\kappa_{\rm loc}x(1+\eta_{\rm loc})^2}{2-q_{\rm loc}}.
\end{align*}
Here \(\eta_{\rm loc}\le3x/4\) and
\(q_{\rm loc}\le x(4+3x)/(3+8x)\). Keep their two errors together:
\begin{align*}
 &\frac{\eta_{\rm loc}^2}{2}
 +\frac{\kappa_{\rm loc}x(1+\eta_{\rm loc})^2}{2-q_{\rm loc}}\\
 &\quad\le\frac9{32}x^2+\frac{x(4+3x)^2}{12(2+4x-x^2)}
 \le\frac23x+\frac15x^2.
\end{align*}
For the last inequality, the identity
\[
 (8-x)(2+4x-x^2)-(4+3x)^2=x(6-21x+x^2)>0
\]
holds for \(x<1/4\). It bounds the rational term by
\(2x/3-x^2/12\); then \(9/32-1/12<1/5\).

For the other entropy, write
\(h_b(p)=\log(1+x^2)+2\ell x^2/(1+x^2)\).
The same trapezoid bound for the logarithm yields
\begin{align*}
 A_{\rm loc}\log(1+x^2)
 &\le\frac12+\frac34x+\frac34x^2\\
 &\quad+\frac{x^3[2x(1+x^2)-3]}{8(1+x^2)}\\
 &\le\frac12+\frac34x+\frac34x^2.
\end{align*}
The remainder is negative for \(x<1/4\). Therefore
\begin{align*}
 A_{\rm loc}h_b(p)\le{}&\frac12+\frac34x+\frac34x^2\\
 &+\frac{2A_{\rm loc}\ell x^2}{1+x^2}.
\end{align*}
The coefficient of \(\ell\) now cancels exactly to
\begin{align*}
 1+\eta_{\rm loc}-\frac{2A_{\rm loc}x^2}{1+x^2}
 &=-x\left[\frac{3}{4(1+x^2)}+x\right].
\end{align*}
Therefore
\begin{equation}
 \begin{split}
 \Lambda_{\rm loc}\ge{}&\frac12-(1+\eta_{\rm loc})\log\kappa_{\rm loc}
 -\frac{17}{12}x-\frac{19}{20}x^2\\
 &-\ell x\left[\frac{3}{4(1+x^2)}+x\right].
 \end{split}
 \label{eq:manual-local-cancellation}
\end{equation}

\paragraph{Uniform remainder estimates.}
The same trapezoid bound at \(t=1/3\) gives
\(\log(4/3)<7/24\). Since \(\eta_{\rm loc}\le3/128\)
and \(\kappa_{\rm loc}<4/3\),
\[
 (1+\eta_{\rm loc})\log\kappa_{\rm loc}
 <\left(1+\frac3{128}\right)\frac7{24}<\frac13.
\]
Also
\[
 \frac{17}{12}x+\frac{19}{20}x^2
 <\frac{17}{12\cdot32}+\frac{19}{20\cdot32^2}<\frac1{20}.
\]
Finally \(\ell e^{-\ell}\) decreases here, so \(\ell x<7/64\), and
\[
 \ell x\left[\frac3{4(1+x^2)}+x\right]
 <\frac7{64}\left(\frac34+\frac1{32}\right)<\frac3{32}.
\]
Substitution in \eqref{eq:manual-local-cancellation} proves
\begin{align*}
 \Lambda_{\rm loc}&>\frac12-\frac13-\frac1{20}-\frac3{32}\\
 &=\frac{11}{480}>0.
\end{align*}
Thus the criterion holds at \(a_{\rm cut}\) throughout
\(\ell\ge7/2\). For fixed \(r\), its left side
\((1-r^2)h(ra)-(1-ar^2)h(r)\) is concave in \(a\) and vanishes
at \(a=1\). It is therefore nonnegative for every
\(a\in[a_{\rm cut},1]\).
